% Inactive bounded Hessian approximation. No bounded nonlinear product is used.
\section{Decreasing admissible approximation}

Let $(X,\omega)$ be a compact Hermitian manifold of dimension $n$, with
$2\leq m<n$, and write
\[
 H_m(w)=(\omega+\ddc w)^m\wedge\omega^{n-m}
\]
for smooth admissible $w$. In this fragment a bounded upper semicontinuous
potential $u$ is admissible if every smooth upper test $q$ at $x$ satisfies
$\omega(x)+\ddc q(x)\in\overline{\Gamma_m(\omega(x))}$.
This is the standard potential-theoretic interpretation of bounded
$(\omega,m)$-subharmonicity. The classical equivalence of viscosity and
distributional subsolutions for a smooth strictly elliptic linear
operator, together with G\aa rding duality, identifies it with the mixed
linear inequalities in the canonical subharmonic representative.
That linear equivalence is a declared input, not a new nonlinear
product theorem.

The representative matters for this extension beyond continuous
potentials: changing the value at one point does not change a current.
For example, the upper semicontinuous function equal to $0$ except for
the value $1$ at one point satisfies every distributional inequality
of the constant $0$. It fails the upper-test condition, since a local
upper test at that point can have arbitrarily negative complex Hessian.
It is therefore excluded by the standard subharmonic definition.

We use two results of Chapter~12: the unique smooth strictly admissible
solution of $H_m(w)=e^{aw+H}\omega^n$ for every $a>0$ and smooth $H$,
and uniform smooth strictly admissible approximation of continuous
admissible potentials. Neither result uses a bounded Hessian product.
Polarized G\aa rding positivity and the classical linear subsolution
characterization remain the algebraic and analytic inputs.

\begin{lemma}\label{draft:607-smooth-majorants}
Every bounded upper semicontinuous real function $u$ on $X$ is the
pointwise decreasing limit of smooth functions $h_j>u$.
\end{lemma}
\begin{proof}
Fix a Riemannian distance $d$ on $X$ and put
\[
 c_j(x)=\sup_{y\in X}\{u(y)-j d(x,y)\}.
\]
Each $c_j$ is Lipschitz, $c_j\geq u$, and $c_{j+1}\leq c_j$.
The supremum is attained. If $y_j$ attains it at a fixed $x$, then
$j d(x,y_j)\leq\sup_X u-u(x)$, so $y_j\to x$.
Upper semicontinuity gives $\limsup_j c_j(x)\leq u(x)$, and hence
$c_j\downarrow u$ pointwise.

Set $\delta_j=4^{-j}$. Smooth approximation of continuous functions
on a compact manifold gives $s_j\in C^\infty(X)$ with
$\|s_j-c_j\|_\infty\leq\delta_j/4$.
Take $h_j=s_j+\delta_j$. Then
\[
 u+\tfrac34\delta_j\leq h_j\leq c_j+\tfrac54\delta_j,
 \qquad
 h_{j+1}\leq c_j+\tfrac5{16}\delta_j
                    < c_j+\tfrac34\delta_j\leq h_j.
\]
The bounds also imply $h_j(x)\to u(x)$ at every point.
\end{proof}

\begin{lemma}\label{draft:607-bounded-penalty-squeeze}
Let $u$ be bounded admissible and $h\geq u$ smooth. Choose
\[
 b>\max\{1,\|H_m(h)/\omega^n\|_\infty\}.
\]
For each positive integer $N$, let $w_N$ be the smooth strictly
admissible solution of
\[
 H_m(w_N)=b e^{N(w_N-h)}\omega^n.
\]
With
\[
 M_N=\|h\|_\infty+\frac{\log b}{N},\qquad
 E_N=\frac{m\log(N+1)+\log b}{N},
\]
one has $\|w_N\|_\infty\leq M_N$ and
\[
 u-E_N-2M_N/N\leq w_N\leq h.
\]
\end{lemma}
\begin{proof}
At a positive maximum of $w_N-h$, the form $\omega+\ddc h$
lies above $\omega+\ddc w_N$ and is admissible there. Monotonicity
would give $b e^{N(w_N-h)}\leq H_m(h)/\omega^n<b$, a contradiction.
Thus $w_N\leq h$. At extrema of $w_N$, comparison with the zero
potential gives
\[
 \min_X h-\frac{\log b}{N}\leq w_N
                \leq\max_X h-\frac{\log b}{N},
\]
so the asserted absolute bound follows.

The upper semicontinuous function $u-(1+N^{-1})w_N$ attains its
maximum at some $x_N$. The smooth function $(1+N^{-1})w_N$ plus
this maximum is an upper test for $u$ there. Put
\[
 T_N=(\omega+\ddc w_N)^{m-1}\wedge\omega^{n-m},
 \qquad d_N=H_m(w_N)/\omega^n>0.
\]
The upper-test cone condition and polarized positivity yield, at $x_N$,
\[
 (N+1)d_N\geq\omega\wedge T_N/\omega^n
                           \geq d_N^{(m-1)/m}.
\]
Consequently $d_N\geq(N+1)^{-m}$ and $w_N-h\geq-E_N$ there.
Using $u\leq h$, the maximum is at most $E_N+M_N/N$.
On all of $X$ this gives
\[
 u\leq(1+N^{-1})w_N+E_N+M_N/N
                     \leq w_N+E_N+2M_N/N.
\]
The argument uses only an upper test for $u$, with no derivative,
continuity assumption or nonlinear product of $u$.
\end{proof}

For a continuous real obstacle $h$, let $P_c(h)$ be the pointwise
supremum of the continuous admissible functions below $h$. Chapter~12
proves that $P_c(h)$ is continuous and admissible. For smooth $h$ this
also follows directly from the preceding estimate applied to every
continuous competitor: $w_N\leq P_c(h)\leq w_N+E_N+2M_N/N$.
The constants are independent of the competitor, and constant
potentials make the family nonempty. This continuity argument precedes
all assertions about an envelope's Hessian measure.

\begin{proposition}\label{draft:607-bounded-envelope}
For every smooth real obstacle $h$, the continuous envelope $P_c(h)$
equals the pointwise supremum of all bounded admissible functions
below $h$. In particular every such $u$ satisfies $u\leq P_c(h)$.
\end{proposition}
\begin{proof}
The solutions $w_N$ of the preceding lemma are continuous admissible
competitors, so $w_N\leq P_c(h)$. For each bounded admissible
$u\leq h$ the same lemma gives
\[
 u\leq w_N+E_N+2M_N/N\leq P_c(h)+E_N+2M_N/N.
\]
Let $N\to\infty$ for this fixed obstacle and fixed $b$.
Thus $P_c(h)$ dominates every bounded competitor. Since $P_c(h)$
itself is a bounded competitor, it equals their supremum.
\end{proof}

\begin{theorem}\label{draft:607-bounded-decreasing-approximation}
Every bounded admissible potential $u$ on $X$ is the pointwise
decreasing limit of smooth strictly admissible potentials.
The sequence may be chosen uniformly bounded.
\end{theorem}
\begin{proof}
Choose smooth $h_j\downarrow u$ by the first lemma and put
$v_j=P_c(h_j)$. Monotonicity of the envelope and the preceding
proposition give
\[
 u\leq v_{j+1}\leq v_j\leq h_j.
\]
Thus $v_j$ are continuous admissible functions decreasing pointwise
to $u$. Set $\varepsilon_j=4^{-j}$. By the continuous approximation
lemma in Chapter~12, choose smooth strictly admissible $a_j$ with
$\|a_j-v_j\|_\infty\leq\varepsilon_j/4$ and put
$u_j=a_j+\varepsilon_j$. Adding a constant preserves strict
admissibility. Moreover
\[
 u+\tfrac34\varepsilon_j\leq u_j,
 \qquad
 u_{j+1}\leq v_j+\tfrac5{16}\varepsilon_j
                <v_j+\tfrac34\varepsilon_j\leq u_j.
\]
The approximation bounds imply $u_j\to u$ pointwise.
Finally, $\inf_X u\leq u_j\leq\sup_X h_1+5/16$, so the sequence
has a common uniform bound.
\end{proof}

This construction supplies decreasing smooth admissible approximants.
It does not by itself identify a limit of their Hessian products.
Replacing a uniform norm in the continuous product proof by pointwise
monotone convergence requires additional capacity, energy or weighted
convergence estimates, with the Hermitian torsion terms retained.
Bounded contact comparison, domination and regularity still depend on
that further argument.
